\documentclass[11pt,a4paper]{article}

\usepackage[margin=2.5cm]{geometry}
\usepackage{graphicx}
\usepackage{amsmath}
\usepackage{hyperref}
\usepackage{booktabs}
\usepackage{enumitem}
\usepackage{titlesec}
\usepackage{caption}
\usepackage{subcaption}
\usepackage{float}

\titleformat{\section}{\large\bfseries}{\thesection}{1em}{}
\titleformat{\subsection}{\normalsize\bfseries}{\thesubsection}{1em}{}

\title{VLM-Driven High-Level Navigation for Industrial AGVs\\
\large Literature Review }

\author{
Pi-Wei Chen \\
Donato Cerciello \\
Lingyu Qiu \\
Juan Camilo España
}

\date{\today}

\begin{document}

\maketitle


\section{LLM-Based Research on AGVs}
\cite{shah2023lm}
presents a system that combines LLMs, VLMs, and visual navigation models to enable long-range navigation from natural language instructions in unmapped environments.

\cite{lu2024optimization} proposes an AGV navigation algorithm based on large model ensemble. By fusing pre-trained LLM, VLM and navigation models, it improves the decision-making flexibility and accuracy in complex industrial environments, and uses knowledge distillation technology to achieve lightweight deployment of visual language models.

\cite{hou2025framework} proposes a task planning framework for automated changeover of a production line based on a large language model (LLM). It represents semantic relationships through a knowledge graph and uses a two-layer planner (high and low layers) to decompose user instructions into skill sequences and codes that can be executed by the robot, thereby realizing the de-management and efficient execution of the production line changeover process.
\cite{lu2025conflict} proposes a BA-A2C path planning method that integrates the A algorithm, behavior cloning (BC), and A2C reinforcement learning for U-shaped automated container terminals. By using the A algorithm to generate expert experience for pre-training, the method overcomes the problem of slow convergence in traditional reinforcement learning and achieves efficient real-time planning and conflict resolution in multi-AGV environments.
\cite{wang2025vision} proposes a mobile robot navigation method for human-centered intelligent manufacturing (HSM). By integrating 3D point cloud reconstruction, VLM-based zero-sample semantic segmentation, and LLM instruction parsing and code generation technology, it effectively solves the challenges of insufficient generalization ability of robots in unstructured industrial environments and environmental perception noise.
\cite{liu2025hierarchical} proposes the first air-ground heterogeneous robot system framework that integrates LLM reasoning and VLM perception. It performs hierarchical task decomposition and global semantic mapping through a large model, and uses a visual language model enhanced by introducing GridMask to achieve high-precision positioning, guiding air and ground robots to collaboratively complete long-term, zero-sample object planning and manipulation tasks.

\cite{zafar2025swarmvlm} proposes a hierarchical multimodal framework that uses LLM for task decomposition and global semantic mapping, and combines VLM enhanced by introducing GridMask to achieve high-precision perception and localization. For the first time, it realizes zero-sample long-term collaborative task execution of air-ground heterogeneous robot systems in dynamic environments.

\cite{tan2025robonavguard} proposes a deep learning solution called RoboNavGuard, which uses the PyraBiNet++ architecture to fuse global and local features to achieve accurate segmentation of complex obstacles (such as plastic bags, dirt, etc.). Combined with the lightweight 3D visual localization model Re3DVG-Small and the new benchmark dataset FragCloud3DRef++, it significantly improves the AGV's ability to understand fragmented point clouds and the navigation accuracy driven by language commands.
 
\cite{huang2025kiterunner} proposes KiteRunner, a language-driven air-ground cooperative navigation strategy. It innovatively combines a global probabilistic graph generated from UAV orthophotos with a diffusion model-driven local path generation, and integrates CLIP and GPT to improve semantic command understanding, significantly optimizing path efficiency and autonomy for long-distance navigation in the open world.

\cite{hu2025integrating} proposes a personalized autonomous driving system integrating GPT-4 is proposed, which utilizes a large model to process visual language navigation (VLN) and knowledge-driven driving decisions, and realizes wireless sharing of local driving experience through roadside units (RSUs). At the execution end, the system combines obstacle repulsion field and DP/QP optimization algorithm to significantly improve navigation completion rate, decision success rate and smoothness of lane change path.


\cite{wang2025visiopath} integrates the zero-shot semantic awareness capabilities of Visual Language Model (VLM) with the rigorous control logic of Model Predictive Control (MPC). It utilizes bird's-eye view (BEV) video processing to generate semantically enhanced initial trajectories and employs elliptical collision avoidance potential fields and differential dynamic programming (DDP) to achieve context-aware decision-making and safe trajectory planning in dynamic and complex traffic environments.

\cite{shah2023lm} propose LM-Nav, a robot navigation system that requires no language-annotated data or model fine-tuning. It organically combines three pre-trained large models through a decoupled architecture: using GPT-3 to parse natural language instructions into landmark sequences, locating these landmarks through visual observation via CLIP, and finally driving the robot with the ViNG navigation model to achieve long-distance, zero-shot instruction following in complex outdoor environments.

\cite{moriconi2025foundation} proposes a navigation system based on the ROS2 behavior tree framework. It utilizes a visual language model (VLM) to extract semantic information from RGB-D data and combines it with a large language model (LLM) to perform context-aware decision-making reasoning based on semantic descriptions. Finally, the behavior tree coordinates robot actions, which significantly improves semantic understanding and decision-making flexibility while ensuring the safety and reliability required in industrial scenarios.





\bibliographystyle{IEEEtran} 
\bibliography{flics}

\end{document}